\documentclass[12pt,a4paper]{ctexbook}
\usepackage{geometry}
\geometry{margin=2.2cm}
\usepackage{setspace}
\setstretch{1.35}
\usepackage{hyperref}
\title{增广贤文}
\author{古典文献整理}
\date{}
\begin{document}
\maketitle
\tableofcontents
\newpage
\section{上集}
昔時賢文，誨汝諄諄。\par
集韻增廣，多見多聞。\par
觀今宜鑑古，無古不成今。\par
知己知彼，將心比心。\par
酒逢知己飲，詩向會人吟。\par
相識滿天下，知心能幾人？\par
相逢好似初相識，到老終無怨恨心。\par
近水知魚性，近山識鳥音。\par
易漲易退山溪水，易反易覆小人心。\par
運去金成鐵，時來鐵似金。\par
讀書須用意，一字值千金。\par
逢人且說三分話，未可全拋一片心。\par
有意栽花花不發，無心插柳柳成蔭。\par
畫虎畫皮難畫骨，知人知面不知心。\par
錢財如糞土，仁義值千金。\par
流水下灘非有意，白雲出岫本無心。\par
當時若不登高望，誰信東流海洋深？\par
路遙知馬力，日久見人心。\par
兩人一般心，無錢堪買金；\par
一人一般心，有錢難買針。\par
相見易得好，久住難爲人。\par
馬行無力皆因瘦，人不風流只爲貧。\par
饒人不是癡漢，癡漢不會饒人。\par
是親不是親，非親卻是親。\par
美不美，鄉中水；親不親，故鄉人。\par
鶯花猶怕春光老，豈可教人枉度春？\par
相逢不飲空歸去，洞口桃花也笑人。\par
紅粉佳人休使老，風流浪子莫教貧。\par
在家不會迎賓客，出門方知少主人。\par
黃芩無假，阿魏無真。\par
客來主不顧，自是無良賓。\par
良賓方不顧，應恐是癡人。\par
貧居鬧市無人問，富在深山有遠親。\par
誰人背後無人說，哪個人前不說人？\par
有錢道真語，無錢語不真。\par
不信但看筵中酒，杯杯先勸有錢人。\par
鬧裏掙錢，靜處安身。\par
來如風雨，去似微塵。\par
長江後浪推前浪，世上新人趕舊人。\par
近水樓臺先得月，向陽花木早逢春。\par
古人不見今時月，今月曾經照古人。\par
先到爲君，後到爲臣。\par
莫道君行早，更有早行人。\par
莫信直中直，須防仁不仁。\par
山中有直樹，世上無直人。\par
自恨枝無葉，莫怨太陽偏。\par
一切都是命，半點不由人。\par
一年之計在於春，一日之計在於寅。\par
一家之計在於和，一生之計在於勤。\par
責人之心責己，恕己之心恕人。\par
守口如瓶，防意如城。\par
寧可人負我，切莫我負人。\par
再三須慎意，第一莫欺心。\par
虎身猶可近，人毒不堪親。\par
來說是非者，便是是非人。\par
遠水難救近火，遠親不如近鄰。\par
有酒有肉多兄弟，急難何曾見一人？\par
人情似紙張張薄，世事如棋局局新。\par
山中也有千年樹，世上難逢百歲人。\par
力微休負重，言輕莫勸人。\par
無錢休入衆，遭難莫尋親。\par
平生不做皺眉事，世上應無切齒人。\par
士者國之寶，儒爲席上珍。\par
若要斷酒法，醒眼看醉人。\par
求人須求大丈夫，濟人須濟急時無。\par
渴時一滴如甘露，醉後添杯不如無。\par
久住令人賤，頻來親也疏。\par
酒中不語真君子，財上分明大丈夫。\par
出家如初，成佛有餘。\par
積金千兩，不如明解經書。\par
養子不教如養驢，養女不教如養豬。\par
有田不耕倉廩虛，有書不讀子孫愚。\par
倉廩虛兮歲月乏，子孫愚兮禮儀疏。\par
聽君一席話，勝讀十年書。\par
人不通今古，馬牛如襟裾。\par
茫茫四海人無數，哪個男兒是丈夫？\par
白酒釀成緣好客，黃金散盡爲收書。\par
救人一命，勝造七級浮屠。\par
城門失火，殃及池魚。\par
庭前生瑞草，好事不如無。\par
欲求生富貴，須下死工夫。\par
百年成之不足，一旦壞之有餘。\par
人心似鐵，官法如爐。\par
善化不足，惡化有餘。\par
水至清則無魚，人太急則無智。\par
知者減半，愚者全無。\par
在家由父，出嫁從夫。\par
癡人畏婦，賢女敬夫。\par
是非終日有，不聽自然無。\par
竹籬茅舍風光好，道院僧房終不如。\par
寧可正而不足，不可邪而有餘。\par
寧可信其有，不可信其無。\par
命裏有時終須有，命裏無時莫強求。\par
道院迎仙客，書堂隱相儒。\par
庭栽棲鳳竹，池養化龍魚。\par
結交須勝己，似我不如無。\par
但看三五日，相見不如初。\par
人情似水分高下，世事如雲任卷舒。\par
會說說都是，不會說無理。\par
磨刀恨不利，刀利傷人指；\par
求財恨不多，財多害自己。\par
知足常足，終身不辱；\par
知止常止，終身不恥。\par
有福傷財，無福傷己。\par
差之毫釐，失之千里。\par
若登高必自卑，若涉遠必自邇。\par
三思而行，再思可矣。\par
動口不如親爲，求人不如求己。\par
小時是兄弟，長大各鄉里。\par
嫉財莫嫉食，怨生莫怨死。\par
人見白頭嗔，我見白頭喜。\par
多少少年郎，不到白頭死。\par
牆有縫，壁有耳。\par
好事不出門，壞事傳千里。\par
若要人不知，除非己莫爲。\par
爲人不做虧心事，半夜敲門心不驚。\par
賊是小人，智過君子。\par
君子固窮，小人窮斯濫矣。\par
富貴多憂，貧窮自在。\par
不以我爲德，反以我爲仇。\par
寧可直中取，不可曲中求。\par
人無遠慮，必有近憂。\par
知我者謂我心憂，不知我者謂我何求？\par
晴天不肯去，直待雨淋頭。\par
成事莫說，覆水難收。\par
是非只爲多開口，煩惱皆因強出頭。\par
忍得一時之氣，免得百日之憂。\par
近來學得烏龜法，得縮頭時且縮頭。\par
懼法朝朝樂，欺公日日憂。\par
人生一世，草長一春。\par
黑髮不知勤學早，轉眼便是白頭翁。\par
月過十五光明少，人到中年萬事休。\par
兒孫自有兒孫福，莫爲兒孫做馬牛。\par
人生不滿百，常懷千歲憂。\par
今朝有酒今朝醉，明日愁來明日憂。\par
路逢險處須迴避，事到臨頭不自由。\par
人貧不語，水平不流。\par
一家養女百家求，一馬不行百馬憂。\par
有花方酌酒，無月不登樓。\par
三杯通大道，一醉解千愁。\par
深山畢竟藏猛虎，大海終須納細流。\par
惜花須檢點，愛月不梳頭。\par
大抵選她肌骨好，不搽紅粉也風流。\par
受恩深處宜先退，得意濃時便可休。\par
莫待是非來入耳，從前恩愛反爲仇。\par
留得五湖明月在，不愁無處下金鉤。\par
休別有魚處，莫戀淺灘頭。\par
去時終須去，再三留不住。\par
忍一句，息一怒，饒一着，退一步。\par
三十不豪，四十不富，五十將來尋死路。\par
生不認魂，死不認屍。\par
一寸光陰一寸金，寸金難買寸光陰。\par
父母恩深終有別，夫妻義重也分離。\par
人生似鳥同林宿，大難來時各自飛。\par
人善被人欺，馬善被人騎。\par
人無橫財不富，馬無夜草不肥。\par
人惡人怕天不怕，人善人欺天不欺。\par
善惡到頭終有報，只盼來早與來遲。\par
黃河尚有澄清日，豈能人無得運時？\par
得寵思辱，居安思危。\par
念念有如臨敵日，心心常似過橋時。\par
英雄行險道，富貴似花枝。\par
人情莫道春光好，只怕秋來有冷時。\par
送君千里，終有一別。\par
但將冷眼觀螃蟹，看你橫行到幾時。\par
見事莫說，問事不知。\par
閒事莫管，無事早歸。\par
假緞染就真紅色，也被旁人說是非。\par
善事可做，惡事莫爲。\par
許人一物，千金不移。\par
龍生龍子，虎生虎兒。\par
龍游淺水遭蝦戲，虎落平原被犬欺。\par
一舉首登龍虎榜，十年身到鳳凰池。\par
十年寒窗無人問，一舉成名天下知。\par
酒債尋常處處有，人生七十古來稀！\par
養兒防老，積穀防饑。\par
雞豚狗彘之畜，無失其時，數口之家，可以無飢矣。\par
當家才知鹽米貴，養子方知父母恩。\par
常將有日思無日，莫把無時當有時。\par
樹欲靜而風不止，子欲養而親不待。\par
時來風送滕王閣，運去雷轟薦福碑。\par
入門休問榮枯事，且看容顏便得知。\par
官清司吏瘦，神靈廟祝肥。\par
息卻雷霆之怒，罷卻虎豹之威。\par
饒人算之本，輸人算之機。\par
好言難得，惡語易施。\par
一言既出，駟馬難追。\par
道吾好者是吾賊，道吾惡者是吾師。\par
路逢俠客須呈劍，不是才人莫獻詩。\par
三人行必有我師焉。\par
擇其善者而從之，其不善者而改之。\par
欲昌和順須爲善，要振家聲在讀書。\par
少壯不努力，老大徒傷悲。\par
人有善願，天必佑之。\par
莫飲卯時酒，昏昏醉到酉。\par
莫罵酉時妻，一夜受孤悽。\par
種麻得麻，種豆得豆。\par
天眼恢恢，疏而不漏。\par
見官莫向前，作客莫在後。\par
寧添一斗，莫添一口。\par
螳螂捕蟬，豈知黃雀在後？\par
不求金玉重重貴，但願兒孫個個賢。\par
一日夫妻，百世姻緣。\par
百世修來同船渡，千世修來共枕眠。\par
殺人一萬，自損三千。\par
傷人一語，利如刀割。\par
枯木逢春猶再發，人無兩度再少年。\par
未晚先投宿，雞鳴早看天。\par
將相頂頭堪走馬，公侯肚內好撐船。\par
富人思來年，窮人想眼前。\par
世上若要人情好，賒去物品莫取錢。\par
生死有命，富貴在天。\par
擊石原有火，不擊乃無煙。\par
人學始知道，不學亦徒然。\par
莫笑他人老，終須還到老。\par
和得鄰里好，猶如拾片寶。\par
但能守本分，終身無煩惱。\par
大家做事尋常，小家做事慌張。\par
大家禮義教子弟，小家兇惡訓兒郎。\par
君子愛財，取之有道。\par
貞婦愛色，納之以禮。\par
善有善報，惡有惡報。\par
不是不報，時候未到。\par
萬惡淫爲首，百行孝當先。\par
人而無信，不知其可也。\par
一人道虛，千人傳實。\par
凡事要好，須問三老。\par
若爭小利，便失大道。\par
家中不和鄰里欺，鄰里不和說是非。\par
年年防飢，夜夜防盜。\par
學者是好，不學不好。\par
學者如禾如稻，不學如草如蒿。\par
遇飲酒時須防醉，得高歌處且高歌。\par
因風吹火，用力不多。\par
不因漁夫引，怎能見波濤？\par
無求到處人情好，不飲任他酒價高。\par
知事少時煩惱少，識人多處是非多。\par
進山不怕傷人虎，只怕人情兩面刀。\par
強中更有強中手，惡人須用惡人磨。\par
會使不在家富豪，風流不用衣着佳。\par
光陰似箭，日月如梭。\par
天時不如地利，地利不如人和。\par
黃金未爲貴，安樂值錢多。\par
爲善最樂，作惡難逃。\par
羊有跪乳之恩，鴉有反哺之情。\par
孝順還生孝順子，忤逆還生忤逆兒。\par
不信但看檐前水，點點滴滴舊窩池。\par
隱惡揚善，執其兩端。\par
妻賢夫禍少，子孝父心寬。\par
已覆之水，收之實難。\par
人生知足時常足，人老偷閒且是閒。\par
處處綠楊堪繫馬，家家有路通長安。\par
既墜釜甑，反顧何益。\par
見者易，學者難。\par
厭靜還思喧，嫌喧又憶山。\par
自從心定後，無處不安然。\par
莫將容易得，便作等閒看。\par
用心計較般般錯，退後思量事事寬。\par
道路各別，養家一般。\par
由儉入奢易，從奢入儉難。\par
知音說與知音聽，不是知音莫與談。\par
點石化爲金，人心猶未足。\par
信了賭，賣了屋。\par
他人觀花，不涉你目。\par
他人碌碌，不涉你足。\par
誰人不愛子孫賢，誰人不愛千鍾粟。\par
奈五行，不是這般題目。\par
莫把真心空計較，兒孫自有兒孫福。\par
書到用時方恨少，事非經過不知難。\par
天下無不是的父母，世上最難得者兄弟。\par
與人不和，勸人養鵝；與人不睦，勸人架屋。\par
但行好事，莫問前程。不交僧道，便是好人。\par
河狹水激，人急計生。\par
明知山有虎，莫向虎山行。\par
路不鏟不平，事不爲不成。\par
無錢方斷酒，臨老始讀經。\par
點塔七層，不如暗處一燈。\par
堂上二老是活佛，何用靈山朝世尊。\par
萬事勸人休瞞昧，舉頭三尺有神明。\par
但存方寸土，留與子孫耕。\par
滅卻心頭火，剔起佛前燈。\par
惺惺多不足，濛濛作公卿。\par
衆星朗朗，不如孤月獨明。\par
兄弟相害，不如友生。\par
合理可作，小利不爭。\par
牡丹花好空入目，棗花雖小結實多。\par
欺老莫欺小，欺人心不明。\par
勤奮耕鋤收地利，他時飽暖謝蒼天。\par
得忍且忍，得耐且耐，不忍不耐，小事成災。\par
相論逞英豪，家計漸漸退。\par
賢婦令夫貴，惡婦令夫敗。\par
一人有慶，兆民鹹賴。\par
人老心未老，人窮志莫窮。\par
人無千日好，花無百日紅。\par
黃蜂一口針，橘子兩邊分。\par
世間痛恨事，最毒淫婦心。\par
殺人可恕，情理不容。\par
乍富不知新受用，乍貧難改舊家風。\par
座上客常滿，杯中酒不空。\par
屋漏更遭連夜雨，行船又遇打頭風。\par
筍因落籜方成竹，魚爲奔波始化龍。\par
記得少年騎竹馬，轉眼又是白頭翁。\par
禮義生於富足，盜賊出於賭博。\par
天上衆星皆拱北，世間無水不朝東。\par
士爲知己者死，女爲悅己者容。\par
色即是空，空即是色。\par
君子安貧，達人知命。\par
良藥苦口利於病，忠言逆耳利於行。\par
順天者昌，逆天者亡。\par
有緣千里來相會，無緣對面不相逢。\par
有福者昌，無福者亡。\par
人爲財死，鳥爲食亡。\par
夫妻相和好，琴瑟與笙簧。\par
紅粉易妝嬌態女，無錢難作好兒郎。\par
有子之人貧不久，無兒無女富不長。\par
善必壽老，惡必早亡。\par
爽口食多偏作病，快心事過恐遭殃。\par
富貴定要依本分，貧窮不必再思量。\par
畫水無風空作浪，繡花雖好不聞香。\par
貪他一斗米，失卻半年糧。\par
爭他一腳豚，反失一肘羊。\par
龍歸晚洞雲猶溼，麝過春山草木香。\par
平生只會說人短，何不回頭把己量？\par
見善如不及，見惡如探湯。\par
人窮志短，馬瘦毛長。\par
自家心裏急，他人未知忙。\par
貧無達士將金贈，病有高人說藥方。\par
觸來莫與競，事過心清涼。\par
秋來滿山多秀色，春來無處不花香。\par
凡人不可貌相，海水不可斗量。\par
清清之水爲土所防，濟濟之士爲酒所傷。\par
蒿草之下或有蘭香，茅茨之屋或有侯王。\par
無限朱門生餓殍，幾多白屋出公卿。\par
酒裏乾坤大，壺中日月長。\par
拂石坐來春衫冷，踏花歸去馬蹄香。\par
萬事前身定，浮生空自忙。\par
叫月子規喉舌冷，宿花蝴蝶夢魂香。\par
一言不中，千言不用。\par
一人傳虛，百人傳實。\par
萬金良藥，不如無疾。\par
千里送鵝毛，禮輕情義重。\par
世事如明鏡，前程暗似漆。\par
君子懷刑，小人懷惠。\par
架上碗兒輪流轉，媳婦自有做婆時。\par
人生一世，如駒過隙。\par
良田萬頃，日食一升。\par
大廈千間，夜眠八尺。\par
千經萬典，孝義爲先。\par
天上人間，方便第一。\par
一字入公門，九牛拔不出。\par
八字衙門向南開，有理無錢莫進來。\par
欲求天下事，須用世間財。\par
富從升合起，貧因不算來。\par
近河不得枉使水，近山不得枉燒柴。\par
家無讀書子，官從何處來？\par
慈不掌兵，義不掌財。\par
一夫當關，萬夫莫開。\par
萬事不由人計較，一生都是命安排。\par
白雲本是無心物，卻被清風引出來。\par
慢行急行，逆取順取。\par
命中只有如許財，絲毫不可有閃失。\par
人間私語，天聞若雷。\par
暗室虧心，神目如電。\par
一毫之惡，勸人莫作。一毫之善，與人方便。\par
虧人是禍，饒人是福，天眼恢恢，報應甚速。\par
聖賢言語，神欽鬼服。\par
人各有心，心各有見。\par
口說不如身逢，耳聞不如目見。\par
見人富貴生歡喜，莫把心頭似火燒。\par
養兵千日，用在一時。\par
國清才子貴，家富小兒嬌。\par
利刀割體瘡猶使，惡語傷人恨不消。\par
公道世間唯白髮，貴人頭上不曾饒。\par
有才堪出衆，無衣懶出門。\par
爲官須作相，及第必爭先。\par
苗從地發，樹由枝分。\par
宅裏燃火，煙氣成雲。\par
以直報怨，知恩報恩。\par
紅顏今日雖欺我，白髮他時不放君。\par
借問酒家何處有，牧童遙指杏花村。\par
父子和而家不退，兄弟和而家不分。\par
一片雲間不相識，三千里外卻逢君。\par
官有公法，民有私約。\par
平時不燒香，臨時抱佛腳。\par
幸生太平無事日，恐防年老不多時。\par
國亂思良將，家貧思良妻。\par
池塘積水須防旱，田地深耕足養家。\par
根深不怕風搖動，樹正何愁月影斜。\par
爭得貓兒，失卻牛腳。\par
愚者千慮，必有一得，智者千慮，必有一失。\par
始吾於人也，聽其言而信其行。\par
今吾於人也，聽其言而觀其行。\par
哪個梳頭無亂髮，情人眼裏出西施。\par
珠沉淵而川媚，玉韞石而山輝。\par
夕陽無限好，只恐不多時。\par
久旱逢甘霖，他鄉遇故知；洞房花燭夜，金榜題名時。\par
惜花春起早，愛月夜眠遲。\par
掬水月在手，弄花香滿衣。\par
桃紅李白薔薇紫，問着東君總不知。\par
教子教孫須教義，栽桑栽柘少栽花。\par
休念故鄉生處好，受恩深處便爲家。\par
學在一人之下，用在萬人之上。\par
一日爲師，終生爲父。\par
忘恩負義，禽獸之徒。\par
勸君莫將油炒菜，留與兒孫夜讀書。\par
書中自有千鍾粟，書中自有顏如玉。\par
莫怨天來莫怨人，五行八字命生成。\par
莫怨自己窮，窮要窮得乾淨；莫羨他人富，富要富得清高。\par
別人騎馬我騎驢，仔細思量我不如，\par
待我回頭看，還有挑腳漢。\par
路上有飢人，家中有剩飯。\par
積德與兒孫，要廣行方便。\par
作善鬼神欽，作惡遭天遣。\par
積錢積穀不如積德，買田買地不如買書。\par
一日春工十日糧，十日春工半年糧。\par
疏懶人沒吃，勤儉糧滿倉。\par
人親財不親，財利要分清。\par
十分伶俐使七分，常留三分與兒孫，\par
若要十分都使盡，遠在兒孫近在身。\par
君子樂得做君子，小人枉自做小人。\par
好學者則庶民之子爲公卿，不好學者則公卿之子爲庶民。\par
惜錢莫教子，護短莫從師。\par
記得舊文章，便是新舉子。\par
人在家中坐，禍從天上落。\par
但求心無愧，不怕有後災。\par
只有和氣去迎人，哪有相打得太平。\par
忠厚自有忠厚報，豪強一定受官刑。\par
人到公門正好修，留些陰德在後頭。\par
爲人何必爭高下，一旦無命萬事休。\par
山高不算高，人心比天高。\par
白水變酒賣，還嫌豬無糟。\par
貧寒休要怨，寶貴不須驕。\par
善惡隨人作，禍福自己招。\par
奉勸君子，各宜守己。\par
只此呈示，萬無一失。\par

\section{下集}
前人俗語，言淺理深。\par
補遺增廣，集成書文。\par
世上無難事，只怕不專心。\par
成人不自在，自在不成人；\par
金憑火煉方知色，與人交財便知心。\par
乞丐無糧，懶惰而成。\par
勤儉爲無價之寶，節糧乃衆妙之門。\par
省事儉用，免得求人。\par
量大禍不在，機深禍亦深。\par
善爲至寶深深用，心作良田世世耕。\par
羣居防口，獨坐防心。\par
體無病爲富貴，身平安莫怨貧。\par
敗家子弟揮金如土，貧家子弟積土成金。\par
富貴非關天地，禍福不是鬼神。\par
安分貧一時，本分終不貧。\par
不拜父母拜乾親，弟兄不和結外人。\par
人過留名，雁過留聲。\par
擇子莫擇父，擇親莫擇鄰。\par
愛妻之心是主，愛子之心是親。\par
事從根起，藕葉連心。\par
禍與福同門，利與害同城。\par
清酒紅人臉，財帛動人心！\par
寧可葷口唸佛，不可素口罵人。\par
有錢能說話，無錢話不靈。\par
豈能盡如人意？但求不愧吾心。\par
不說自己井繩短，反說他人箍井深。\par
恩愛多生病，無錢便覺貧。\par
只學斟酒意，莫學下棋心。\par
孝莫假意，轉眼便爲人父母。\par
善休望報，回頭只看汝兒孫！\par
口開神氣散，舌出是非生！\par
彈琴費指甲，說話費精神。\par
千貫買田，萬貫結鄰。\par
人言未必猶盡，聽話只聽三分。\par
隔壁豈無耳，窗外豈無人？\par
財可養生須注意，事不關己不勞心。\par
酒不護賢，色不護病；\par
財不護親，氣不護命！\par
一日不可無常業，安閒便易起邪心！\par
炎涼世態，富貴更甚於貧賤；\par
嫉妒人心，骨肉更甚於外人！\par
瓜熟蒂落，水到渠成。\par
人情送匹馬，買賣不饒針！\par
過頭飯好吃，過頭話難聽！\par
事多累了自己，田多養了衆人。\par
怕事忍事不生事自然無事；\par
平心靜心不欺心何等放心！\par
天子至尊不過於理，在理良心天下通行。\par
好話不在多說，有理不在高聲！\par
一朝權在手，便把令來行。\par
甘草味甜人可食，巧言妄語不可聽。\par
當場不論，過後枉然。\par
貧莫與富鬥，富莫與官爭！\par
官清難逃猾吏手，衙門少有唸佛人！\par
家有千口，主事一人。\par
父子竭力山成玉，弟兄同心土變金。\par
當事者迷，旁觀者清。\par
怪人不知理，知理不怪人。\par
未富先富終不富，未貧先貧終不貧。\par
少當少取，少輸當贏！\par
飽暖思淫慾，飢寒起盜心！\par
蚊蟲遭扇打，只因嘴傷人！\par
欲多傷神，財多累心！\par
布衣得暖真爲福，千金平安即是春。\par
家貧出孝子，國亂顯忠臣！\par
寧做太平犬，莫做離亂人！\par
人有幾等，官有幾品。\par
理不衛親，法不爲民。\par
自重者然後人重，人輕者便是自輕。\par
自身不謹，擾亂四鄰。\par
快意事過非快意，自古敗名因敗事。\par
傷身事莫做，傷心話莫說。\par
小人肥口，君子肥身。\par
地不生無名之輩，天不生無路之人。\par
一苗露水一苗草，一朝天子一朝臣。\par
讀未見書，如得良友；見已讀書，如逢故人。\par
福滿須防有禍，兇多料必無爭。\par
不怕三十而死，只怕死後無名。\par
但知江湖者，都是薄命人。\par
不怕方中打死人，只知方中無好人。\par
說長說短，寧說人長莫說短；\par
施恩施怨，寧施人恩莫施怨。\par
育林養虎，虎大傷人。\par
冤家抱頭死，事要解交人。\par
捲簾歸乳燕，開扇出蒼蠅。\par
愛鼠常留飯，憐蛾燈罩紗。\par
人命在天，物命在人。\par
奸不通父母，賊不通地鄰。\par
盜賊多出賭博，人命常出姦情。\par
治國信讒必殺忠臣，治家信讒必疏其親。\par
治國不用佞臣，治家不用佞婦。\par
好臣一國之寶，好婦一家之珍。\par
穩的不滾，滾的不穩。\par
兒不嫌母醜，狗不嫌家貧。\par
君子千錢不計較，小人一錢惱人心。\par
人前顯貴，鬧裏奪爭。\par
要知江湖深，一個不做聲。\par
知止自當出妄想，安貧須是禁奢心。\par
初入行業，三年事成；\par
初吃饅頭，三年口生。\par
家無生活計，坐吃如山崩。\par
家有良田萬頃，不如薄藝在身；\par
藝多不養家，食多嚼不贏。\par
命中只有八合米，走遍天下不滿升。\par
使心用心，反害自身。\par
國家無空地，世上無閒人。\par
妙藥難醫怨逆病，混財不富窮命人。\par
耽誤一年春，十年補不清；\par
人能處處能，草能處處生。\par
會打三班鼓，也要幾個人。\par
人不走不親，水不打不渾。\par
三貧三富不到老，十年興敗多少人！\par
買貨買得真，折本折得輕；\par
不怕問到，只怕倒問。\par
人強不如貨強，價高不如口便。\par
會買買怕人，會賣賣怕人。\par
只只船上有梢公，天子足下有貧親。\par
既知莫望，不知莫向。\par
在一行，練一行；\par
窮莫失志，富莫癲狂。\par
天欲令其滅亡，必先讓其瘋狂。\par
梢長人膽大，梢短人心慌。\par
隔行莫貪利，久煉必成鋼。\par
瓶花雖好豔，相看不耐長。\par
早起三光，遲起三慌。\par
未來休指望，過去莫思量；\par
時來遇好友，病去遇良方。\par
布得春風有夏雨，哈得秋風大家涼。\par
晴帶雨傘，飽帶飢糧。\par
滿壺全不響，半壺響叮噹。\par
久利之事莫爲，衆爭之地莫往。\par
老醫迷舊疾，朽藥誤良方；\par
該在水中死，不在岸上亡。\par
舍財不如少取，施藥不如傳方。\par
倒了城牆醜了縣官，打了梅香醜了姑娘。\par
燕子不進愁門，耗子不鑽空倉。\par
蒼蠅不叮無縫蛋，謠言不找謹慎人。\par
一人舍死，萬人難當。\par
人爭一口氣，佛爭一炷香。\par
門爲小人而設，鎖乃君子之防。\par
舌咬只爲揉，齒落皆因眶。\par
硬弩弦先斷，鋼刀刃自傷。\par
賊名難受，龜名難當。\par
好事他人未見講，錯處他偏說得長。\par
男子無志純鐵無鋼，女子無志爛草無瓤。\par
生男欲得成龍猶恐成獐，生女欲得成鳳猶恐成虎。\par
養男莫聽狂言，養女莫叫離母。\par
男子失教必愚頑，女子失教定粗魯。\par
生男莫教弓與弩，生女莫教歌與舞。\par
學成弓弩沙場災，學成歌舞爲人妾。\par
財交者密，財盡者疏。\par
婚姻論財，夫妻之道。\par
色嬌者親，色衰者疏。\par
少實勝虛，巧不如拙。\par
百戰百勝不如無爭，萬言萬中不如一默。\par
有錢不置怨逆產，冤家宜解不宜結。\par
近朱者赤，近墨者黑。\par
一個山頭一隻虎，惡龍難鬥地頭蛇。\par
出門看天色，進門看臉色。\par
商賈買賣如施捨，買賣公平如積德。\par
天生一人，地生一穴。\par
家無三年之積不成其家，國無九年之積不成其國。\par
男子有德便是才，女子無才便是德。\par
有錢難買子孫賢，女兒不請上門客。\par
男大當婚女大當嫁，不婚不嫁惹出笑話。\par
謙虛美德，過謙即詐。\par
自己跌倒自己爬，望人扶持都是假。\par
人不知己過，牛不知力大。\par
一家飽暖千家怨，一物不見賴千家。\par
當面論人惹恨最大，是與不是隨他說吧！\par
誰人做得千年主，轉眼流傳八百家。\par
滿載芝麻都漏了，還在水裏撈油花！\par
皇帝坐北京，以理統天下。\par
五百年前共一家，不同祖宗也同華！\par
學堂大如官廳，人情大過王法。\par
找錢猶如針挑土，用錢猶如水推沙！\par
害人之心不可有，防人之心不可無！\par
不愁無路，就怕不做。\par
須向根頭尋活計，莫從體面下功夫！\par
禍從口出，病從口入。\par
藥補不如肉補，肉補不如養補。\par
思慮之害甚於酒色，日日勞力上牀呼疾。\par
人怕不是福，人欺不是辱。\par
能言不是真君子，善處方爲大丈夫！\par
爲人莫犯法，犯法身無主。\par
姊妹同肝膽，弟兄同骨肉。\par
慈母多誤子，悍婦必欺夫！\par
君子千里同舟，小人隔牆易宿。\par
文錢逼死英雄漢，財不歸身恰是無。\par
妻子如衣服，弟兄似手足。\par
衣服補易新，手足斷難續。\par
盜賊怨失主，不孝怨父母。\par
一時勸人以口，百世勸人以書。\par
我不如人我無其福，人不如我我常知足！\par
撿金不忘失金人，三兩黃銅四兩福。\par
因禍得福，求賭必輸。\par
一言而讓他人之禍，一忿而折平生之福。\par
天有不測風雲，人有旦夕禍福。\par
不淫當齋，淡飽當肉。\par
緩步當車，無禍當福。\par
男無良友不知己之有過，女無明鏡不知面之精粗。\par
事非親做，不知難處。\par
十年易讀舉子，百年難淘江湖！\par
積錢不如積德，閒坐不如看書。\par
思量挑擔苦，空手做是福。\par
時來易借銀千兩，運去難賒酒半壺。\par
天晴打過落雨鋪，少時享過老來福。\par
與人方便自己方便，一家打牆兩家好看。\par
當面留一線，過後好相見。\par
入門掠虎易，開口告人難。\par
手指要往內撇，家醜不可外傳。\par
浪子出於祖無德，孝子出於前人賢。\par
貨離鄉貴，人離鄉賤。\par
樹挪死，人挪活。\par
在家千日好，出門處處難。\par
三員長者當官員，幾個明人當知縣？\par
明人自斷，愚人官斷。\par
人怕三見面，樹怕一墨線。\par
村夫硬似鐵，光棍軟如棉。\par
不是撐船手，怎敢拿篙竿！\par
天下禮儀無窮，一人知識有限。\par
一人不得二人計，宋江難結萬人緣。\par
家有三畝田，不離衙門前，鄉間無強漢，衙門就餓飯。\par
人人依禮儀，天下不設官。\par
衙門錢，眼睛錢；\par
田禾錢，千萬年。\par
詩書必讀，不可做官。\par
爲人莫當官，當官皆一般。\par
換了你我去，恐比他還貪。\par
官吏清廉如修行，書差方便如行善。\par
靠山吃山，種田吃田。\par
吃盡美味還是鹽，穿盡綾羅還是棉。\par
一夫不耕，全家餓飯，一女不織，全家受寒。\par
金銀到手非容易，用時方知來時難。\par
先講斷，後不亂，免得藕斷絲不斷。\par
聽人勸，得一半。\par
不怕慢，只怕站。\par
逢快莫趕，逢賤莫懶。\par
謀事在人，成事在天！\par
長路人挑擔，短路人賺錢。\par
寧賣現二，莫賣賒三。\par
賺錢往前算，折本往後算。\par
小小生意賺大錢，七十二行出狀元。\par
自己無運至，卻怨世界難。\par
膽大不如膽小，心寬甚如屋寬。\par
妻賢何愁家不富，子孫何須受祖田。\par
是兒不死，是財不散。\par
財來生我易，我去生財難。\par
十月灘頭坐，一日下九灘。\par
結交一人難上難，得罪一人一時間。\par
借債經商，賣田還債；\par
賒錢起屋，賣屋還錢。\par
修起廟來鬼都老，拾得秤來姜賣完。\par
不嫖莫轉，不賭莫看。\par
節食以去病，少食以延年。\par
豆腐多了是包水，梢公多了打爛船。\par
無口過是，無眼過難。\par
無身過易，無心過難。\par
不會鳧水怨河灣，不會犁田怨枷擔。\par
他馬莫騎，他弓莫挽。\par
要知心腹事，但聽口中言。\par
寧在人前全不會，莫在人前會不全。\par
事非親見，切莫亂談。\par
打人莫打臉，罵人莫罵短。\par
好言一句三冬暖，話不投機六月寒。\par
人上十口難盤，帳上萬元難還。\par
放債如施，收債如討。\par
告狀討錢，海底摸鹽。\par
衙門深似海，弊病大如天。\par
銀錢莫欺騙，牛馬不好變。\par
好漢莫被人識破，看破不值半文錢。\par
狗咬對頭人，雷打三世冤。\par
不賣香燒無剩錢，井水不打不滿邊。\par
事寬則園，太久則偏。\par
高人求低易，低人求高難。\par
有錢就是男子漢，無錢就是漢子難。\par
人上一百，手藝齊全。\par
難者不會，會者不難。\par
生就木頭造就船，砍的沒得車的圓。\par
心不得滿，事不得全。\par
鳥飛不盡，話說不完。\par
人無喜色休開店，事不遂心莫怨天。\par
選婿莫選田園，選女莫選嫁奩。\par
紅顏女子多薄命，福人出在醜人邊。\par
人將禮義爲先，樹將花果爲園。\par
臨危許行善，過後心又變。\par
天意違可以人回，命早定可以心挽。\par
強盜口內出赦書，君子口中無戲言。\par
貴人語少，貧子話多。\par
快裏須斟酌，耽誤莫遲春。\par
讀過古華佗，不如見症多。\par
東屋未補西屋破，前帳未還後又拖。\par
今年又說明年富，待到明年差不多。\par
志不同己，不必強合。\par
莫道坐中安樂少，須知世上苦情多。\par
本少利微強如坐，屋檐水也滴得多。\par
勤儉持家富，謙恭受益多。\par
細處不斷粗處斷，黃梅不落青梅落。\par
見錢起意便是賊，順手牽羊乃爲盜。\par
要做快活人，切莫尋煩惱。\par
要做長壽人，莫做短命事。\par
要做有後人，莫做無後事。\par
不經一事，不長一智。\par
寧可無錢使，不可無行止。\par
栽樹要栽松柏，結交要結君子。\par
秀才不出門，能知天下事。\par
錢多不經用，兒多不耐死。\par
弟兄爭財家不窮不止，妻妾爭風夫不死不止。\par
男人有志，婦人有勢。\par
夫人死百將臨門，將軍死一卒不至。\par
天旱誤甲子，人窮誤口齒。\par
百歲無多日，光陰能幾時？\par
父母養其身，自己立其志。\par
待有餘而濟人，終無濟人之日；\par
待有閒而讀書，終無讀書之時。\par
此書傳後世，句句必精讀，其中禮和義，奉勸告世人。\par
勤奮讀，苦發奮，走遍天涯如遊刃。\par

\end{document}